\documentclass[11pt]{article}
\usepackage{amsmath,amssymb,amsthm,fullpage}
\newtheorem{theorem}{Theorem}
\newtheorem{proposition}[theorem]{Proposition}
\newtheorem{lemma}[theorem]{Lemma}
\newtheorem{corollary}[theorem]{Corollary}
\theoremstyle{definition}
\newtheorem{remark}[theorem]{Remark}
\DeclareMathOperator{\nuu}{\nu}

\title{Structural results on the factorial divisibility $((n+k)!)^2 \mid (2n)!$\\
{\large (Erd\H{o}s problem 727): a digit trichotomy, a membership engine,\\ and a
first-moment obstruction}}
\author{Autonomous proof-search run (Claude Code)}
\date{28 July 2026}

\begin{document}
\maketitle

\begin{abstract}
For fixed $k\ge 1$ let $S_k=\{n\ge 1: ((n+k)!)^2 \mid (2n)!\}$. Erd\H{o}s--Graham--Ruzsa--Straus
asked whether $S_k$ is infinite for every fixed $k\ge2$; this is open even for $k=2$
(Balakran proved $S_1$ infinite). We do \emph{not} resolve the question. We record four
things established rigorously during an autonomous search, each machine-verified:
(i) a normal form placing 727 and the recently resolved problems 728/729/401 in the same
shape, together with a \emph{digit trichotomy} that explains exactly why the carry-engineering
method that settled the latter cannot settle the former (Theorem~\ref{thm:mirror});
(ii) a corollary deriving the necessary $\sqrt{2n}$-smoothness of the window
$n+1,\dots,n+k$ from the trichotomy rather than assuming it (Corollary~\ref{cor:smooth});
(iii) a membership engine reducing $n \in S_k$ to a cofactor congruence at each large prime
(Theorem~\ref{thm:engine}); and (iv) a quantitative obstruction showing that no unconditional
first-moment argument can succeed at any smoothness threshold (Proposition~\ref{prop:fm}).
Together these localize the difficulty of 727 in a single joint statement about smoothness of
consecutive integers together with cofactor congruences.
\end{abstract}

\section{Notation and the criterion}

Throughout $p$ denotes a prime, $\nu_p$ the $p$-adic valuation, $s_p(x)$ the sum of base-$p$
digits of $x$, and
\[
  \kappa_p(x) := \#\{\text{carries when adding } x+x \text{ in base } p\} = \nu_p\!\binom{2x}{x},
\]
the last equality being Kummer's theorem. We write $m=n+k$.

\begin{lemma}[Criterion]\label{lem:crit}
For all $n \ge k \ge 1$ the following are equivalent:
\begin{enumerate}
\item[(a)] $n \in S_k$;
\item[(b)] $2s_p(n+k)-s_p(2n)\ge 2k$ for every prime $p$;
\item[(c)] $\kappa_p(n) \ge 2\sum_{j=1}^{k}\nu_p(n+j)$ for every prime $p$;
\item[(d)] $\prod_{i=0}^{2k-1}(2m-i) \;\Big|\; \binom{2m}{m}$, i.e.\
      $\kappa_p(m) \ge \nu_p\!\big(\prod_{i=0}^{2k-1}(2m-i)\big)$ for every $p$.
\end{enumerate}
\end{lemma}

\begin{proof}
By Legendre, $\nu_p(x!)=(x-s_p(x))/(p-1)$, so $\nu_p((2n)!)\ge 2\nu_p((n+k)!)$ is
$2n-s_p(2n)\ge 2(n+k)-2s_p(n+k)$, which is (b). For (c), use
$(n+k)!=n!\prod_{j=1}^k(n+j)$ to get
$\nu_p((2n)!)-2\nu_p((n+k)!)=\nu_p\binom{2n}{n}-2\sum_{j\le k}\nu_p(n+j)$ and Kummer.
For (d), $(2m)! = (2n)!\prod_{i=0}^{2k-1}(2m-i)$ and $m!=(n+k)!$ give
$\nu_p((2n)!)-2\nu_p((n+k)!) = \kappa_p(m)-\nu_p\big(\prod_{i<2k}(2m-i)\big)$.
\end{proof}

\noindent\emph{Verification.} All four forms were checked to agree exhaustively for
$k\le 4$, $n<400$ against exact big-integer division, and on $2\cdot10^4$ random triples
(\texttt{experiments/verify\_identities.py}, \texttt{experiments/mirror\_theorem.py}, claim M1).

\section{The digit trichotomy and the mirror}

Form (d) of Lemma~\ref{lem:crit} exhibits 727 as: \emph{does a product of $2k$ consecutive
integers ending at $2m$ divide the central binomial coefficient $\binom{2m}{m}$?} The
problems 728/729/401, in the treatment of \cite{Sot26}, reduce to: \emph{does
$(m+1)\cdots(m+k)$ divide $\binom{2m}{m}$?} (up to the slack $\nu_p(k!)$). The supply side
$\kappa_p(m)$ is identical in both; only the position of the divisor block relative to $m$
differs. The following trichotomy shows that this difference is decisive.

\begin{lemma}[Trichotomy]\label{lem:tri}
Let $p>2k$ be prime and $J\ge1$. Consider the base-$p$ digits of $m$ at positions
$0,\dots,J-1$ and the carries they produce when doubling $m$.
\begin{enumerate}
\item[\textup{(A)}] If $m\equiv -i \pmod{p^J}$ with $1\le i\le k$, those digits are
  $p-i,p-1,\dots,p-1$, all $\ge\lceil p/2\rceil$; they produce $J$ carries.
\item[\textup{(B)}] If $p$ is odd and $2m\equiv i \pmod{p^J}$ with $i$ odd, $0<i<2k$, those
  digits are $\tfrac{p+i}{2},\tfrac{p-1}{2},\dots,\tfrac{p-1}{2}$; they produce $J$ carries.
\item[\textup{(C)}] If $m\equiv +i \pmod{p^J}$ with $0\le i\le k-1$, those digits are
  $i,0,\dots,0$; they produce \emph{no} carry, and no carry enters position $J$.
\end{enumerate}
\end{lemma}

\begin{proof}
(A) A digit $d\ge \lceil p/2\rceil$ satisfies $2d\ge p$, so it carries irrespective of the
incoming carry. (B) At position $0$ we have $2\cdot\frac{p+i}{2}=p+i\ge p$, a carry; at each
subsequent position $2\cdot\frac{p-1}{2}+1=p$, again a carry. (C) At position $0$,
$2i<2k<p$, no carry; at positions $1,\dots,J-1$ the digit is $0$ and the incoming carry is
$0$, so $0$ again. Hence no carry is produced below position $J$ and none enters position $J$.
\end{proof}

\begin{theorem}[Mirror Theorem]\label{thm:mirror}
Let $p>2k$. In the divisor block of the $728/729/401$ normal form, $\{m+1,\dots,m+k\}$, every
prime power divides a term of type \textup{(A)}, so all required carries are free. In the
divisor block of $727$, $\{2m,2m-1,\dots,2m-2k+1\}$, the odd terms are of type \textup{(B)}
(free), while the even terms $2m-2i=2(m-i)$, $0\le i\le k-1$, are of type \textup{(C)}: for
$p^J\|m-i$ the low $J$ digits of $m$ supply no carries at all, and all $J$ carries required by
Lemma~\ref{lem:crit}(d) must come from digit positions $\ge J$.
\end{theorem}

\begin{proof}
For $p>2k$ at most one term of either block is divisible by $p$, since the terms differ by
less than $2k<p$. For the block $\{m+i\}_{i=1}^k$, $p^J\mid m+i$ is case (A). For 727's block,
$\nu_p(2m-i)$ with $i$ odd is case (B), and for $i=2i'$ even, $\nu_p(2(m-i'))=\nu_p(m-i')$
for odd $p$, which is case (C). The carry count $\kappa_p(m)$ decomposes as the carries below
position $J$ plus those at positions $\ge J$; by Lemma~\ref{lem:tri}(C) the former is $0$.
\end{proof}

\noindent\emph{Verification.} Claims M2--M4 of \texttt{experiments/mirror\_theorem.py} tested
(A), (B), (C) on $5763$, $4626$ and $5778$ structured instances respectively, with zero
violations.

\begin{corollary}[Smoothness is forced]\label{cor:smooth}
Let $n>2k^2$ and $n\in S_k$. Then every prime power $p^J\|n+j$ with $p>2k$, $1\le j\le k$,
satisfies $p^{2J}\le p\,m$; in particular $n+1,\dots,n+k$ are all $\sqrt{2n}$-smooth up to the
primes $\le 2k$.
\end{corollary}

\begin{proof}
Such a $p^J$ makes $2(n+j)=2m-2(k-j)$ a type-(C) term, so by Theorem~\ref{thm:mirror} all $J$
required carries occur at positions $\ge J$. There are at most $\lfloor \log_p m\rfloor-J+1$
such positions, whence $J\le \log_p m-J+1$.
\end{proof}

\noindent\emph{Verification.} M5: no $n<4\cdot10^4$ with $k\in\{2,3\}$ violates this.

\section{A membership engine}

Theorem~\ref{thm:mirror} shows the residual content of 727 is a condition on \emph{cofactors}.

\begin{theorem}[Cofactor engine]\label{thm:engine}
Fix $k\ge2$ and $P_0\ge 2k$. Let $n$ satisfy
\begin{enumerate}
\item[\textup{(i)}] $\kappa_\ell(n)\ \ge\ 2\sum_{j=1}^{k}\nu_\ell(n+j)$ for every prime
  $\ell\le P_0$; and
\item[\textup{(ii)}] for every $j\le k$ and every prime $\ell>P_0$ with $\ell\mid n+j$:
  $\ell\,\|\,n+j$ and, writing $W=(n+j)/\ell$,
  \[ (W-1)\bmod \ell \ \ge\ \tfrac{\ell-1}{2}. \]
\end{enumerate}
Then $n\in S_k$.
\end{theorem}

\begin{proof}
By Lemma~\ref{lem:crit}(c) it suffices to check $\kappa_\ell(n)\ge2\sum_j\nu_\ell(n+j)$ at
every prime. For $\ell\le P_0$ this is (i). For $\ell>P_0\ge 2k$ at most one window element is
divisible by $\ell$ (the elements differ by less than $k<\ell$), and by (ii) that valuation is
$1$, so the demand is $2$. Write $n=\ell W-j$. The two lowest base-$\ell$ digits of $n$ are
$\ell-j$ and $(W-1)\bmod \ell$. Doubling: position $0$ gives $2(\ell-j)\ge\ell$ since
$\ell>2k\ge 2j$, a carry; position $1$ receives that carry and gives
$2\big((W-1)\bmod\ell\big)+1\ge\ell$ by (ii), a second carry. Hence $\kappa_\ell(n)\ge2$.
\end{proof}

\begin{remark}
Condition (ii) is sufficient but not necessary (carries at higher positions also count):
over all even $n\le 6\cdot10^4$ the predicate of Theorem~\ref{thm:engine} has zero false
positives and certifies $45$ of the $263$ members of $S_2$; for $k=3$ it certifies $10$ of the
$41$ members of $S_3\cap[1,6\cdot10^4]$, including the least member $3475$. A prime-parametrized
specialization (each window element a balanced semiprime in an explicit ratio box) was verified
on $123$ solutions and independently regenerated.
\end{remark}

\section{A first-moment obstruction}

By Theorem~\ref{thm:mirror}, for $\ell\|n+j$ the condition at $\ell$ is exactly
$\kappa_\ell(W)\ge1$ with $W=(n+j)/\ell$, i.e.\ \emph{the cofactor $W$ has some base-$\ell$
digit $\ge\lceil \ell/2\rceil$}. Failure is the event that $W$ is ``digit-poor'' base $\ell$,
of density $\asymp 2^{-D}$ where $D\approx \log n/\log \ell$ is the number of usable digits.
This makes the following accounting unavoidable. Let
\[
 F(u) := \mathrm{dens}\{n:\ \exists\, \ell\le n^{1/u},\ \ell\|(n+j),\ W \text{ digit-poor}\},
 \qquad
 S(u) := \mathrm{dens}\{n:\ n+1,n+2 \text{ both } n^{1/u}\text{-smooth}\}.
\]

\begin{proposition}[First-moment obstruction]\label{prop:fm}
$F(u)\asymp 2^{-u}$ while $S(u)\sim\rho(u)^2\approx u^{-2u}$, so $S(u)/F(u)\to0$; consequently
no unconditional first-moment or union-bound argument over all $n$ can produce a member of
$S_2$, at any smoothness threshold.
\end{proposition}

\noindent Measured exhaustively for $n\le2\cdot10^5$
(\texttt{experiments/first\_moment\_obstruction.py}):
\[
\begin{array}{c|ccccc}
u & 2.0 & 2.5 & 3.0 & 4.0 & 5.0\\\hline
F(u) & 0.4096 & 0.2581 & 0.1685 & 0.0793 & 0.0327\\
2^{-u} & 0.2500 & 0.1768 & 0.1250 & 0.0625 & 0.0313\\
S(u) & 0.0592 & 0.0091 & 0.0009 & 0.00002 & \approx0\\
S/F & 0.145 & 0.035 & 0.006 & 0.0003 & \approx0
\end{array}
\]
At $u=2$, $98.96\%$ of all $n\le3\cdot10^5$ fail some condition, against a smooth-window
density of $9.37\%\approx\rho(2)^2$. Moreover the conditions are positively correlated in
failure: the observed conditional membership rate is $0.1105$ versus $0.3297$ predicted by an
independence model.

\section{What remains}

Combining Theorem~\ref{thm:engine} with a proved small-prime counting lemma (density
$1-18\exp(-\sqrt{\log M}/120)$ for the exact criterion at all $p\le\exp(c\sqrt{\log M})$,
uniformly in arithmetic progressions of modulus $\le M^{1/10}$), the problem for fixed $k$
reduces to a single joint statement: \emph{infinitely many $n$ whose window
$n+1,\dots,n+k$ is $\sqrt{2n}$-smooth and whose large prime factors all satisfy the cofactor
congruence of Theorem~\ref{thm:engine}\textup{(ii)}}. For $k=2$ the smooth supply alone is
known at positive density (Hildebrand's theorem on Balog's conjecture); by
Proposition~\ref{prop:fm} the congruences cannot be removed by an unconditional first moment,
so they must be counted inside the smooth set. For $k\ge3$ the supply itself is open:
positive-density $k$-strings of $n^{\alpha}$-smooth integers are known only for
$\alpha>e^{-1/(k-1)}$, which exceeds $1/2$ precisely when $k\ge3$.

\begin{thebibliography}{9}
\bibitem{Bal29} M.~Balakran, \emph{On the values of $n$ for which $(n+1)!^2$ divides $(2n)!$},
  1929.
\bibitem{EGRS75} P.~Erd\H{o}s, R.~L.~Graham, I.~Z.~Ruzsa, E.~G.~Straus,
  \emph{On the prime factors of $\binom{2n}{n}$}, Math.\ Comp.\ \textbf{29} (1975), 83--92.
\bibitem{Erd68} P.~Erd\H{o}s, \emph{Aufgabe 557}, Elem.\ Math.\ \textbf{23} (1968), 111--113.
\bibitem{Hil85} A.~Hildebrand, \emph{On a conjecture of Balog}, Proc.\ Amer.\ Math.\ Soc.\
  \textbf{95} (1985), 517--523.
\bibitem{BW98} A.~Balog, T.~D.~Wooley, \emph{On strings of consecutive integers with no large
  prime factors}, J.\ Austral.\ Math.\ Soc.\ Ser.\ A \textbf{64} (1998), 266--276.
\bibitem{Pom15} C.~Pomerance, \emph{Divisors of the middle binomial coefficient},
  Amer.\ Math.\ Monthly \textbf{122} (2015), 636--644.
\bibitem{Sot26} N.~Sothanaphan, \emph{Resolution of Erd\H{o}s Problem \#728: a writeup of
  Aristotle's Lean proof}, arXiv:2601.07421 (2026). \emph{Used here only for the normal form
  and method background; no result about problem 727 is imported.}
\end{thebibliography}

\end{document}
